%% ============================================================
%% CodeChangeHistory.tex — ARES 코드 수정 이력 (롤링 보고서)
%% 2026-05-19 ~ 진행 중
%% BLE 파이프라인, 부팅 안정화, DC 모터 대칭, BATCH, FinalBuild,
%% 12차시 × 4미션 웹 재편, 인라인 자산 shim 오프라인 빌드
%% XeLaTeX 컴파일: xelatex -shell-escape CodeChangeHistory.tex
%% ============================================================
\documentclass[11pt, oneside]{article}

%% ----- 한국어 / 폰트 -----
\usepackage{kotex}
\usepackage{fontspec}
\defaultfontfeatures{Mapping=}
%% 사용 환경에 맞게 수정. macOS: /Users/<USER>/Library/Fonts/
\newcommand{\userfontpath}{C:/Users/young/AppData/Local/Microsoft/Windows/Fonts/}

\setmainfont{KoPub Batang}[
  Path=\userfontpath, Scale=0.92, Mapping=, LetterSpace=-2.0,
  UprightFont = KoPub Batang Medium.ttf,
  BoldFont    = KoPub Batang Bold.ttf
]
\setsansfont{KoPub Dotum}[
  Path=\userfontpath, Scale=0.92, Mapping=, LetterSpace=-2.0,
  UprightFont = KoPub Dotum Medium.ttf,
  BoldFont    = KoPub Dotum Bold.ttf
]
\setmonofont{D2Coding}[
  Path=\userfontpath, Scale=0.88,
  Extension = .ttc,
  UprightFont = D2Coding-Ver1.3.2-20180524-all
]
\setmainhangulfont{KoPub Batang}[
  Path=\userfontpath, Scale=0.92, Mapping=, LetterSpace=-2.0,
  UprightFont = KoPub Batang Medium.ttf,
  BoldFont    = KoPub Batang Bold.ttf
]
\setsanshangulfont{KoPub Dotum}[
  Path=\userfontpath, Scale=0.92, Mapping=, LetterSpace=-2.0,
  UprightFont = KoPub Dotum Medium.ttf,
  BoldFont    = KoPub Dotum Bold.ttf
]
\setmonohangulfont{KoPub Batang}[
  Path=\userfontpath, Scale=0.92, LetterSpace=-2.0,
  UprightFont = KoPub Batang Medium.ttf,
  BoldFont    = KoPub Batang Bold.ttf
]

%% ----- 레이아웃 -----
\usepackage{geometry}
\geometry{
  paperwidth=190mm, paperheight=240mm,
  top=24mm, bottom=28mm, inner=24mm, outer=22mm,
  headheight=14pt, headsep=8mm, footskip=28pt
}
\linespread{1.12}
\setlength{\parskip}{0.85em}
\setlength{\parindent}{0pt}
\raggedbottom
\setcounter{tocdepth}{2}
\setcounter{secnumdepth}{2}

%% ----- 색상 / 박스 / 코드 -----
\usepackage{xcolor}
\definecolor{deepblue}{RGB}{22, 56, 100}
\definecolor{accentcopper}{RGB}{200, 105, 50}
\definecolor{labcolor}{RGB}{40, 140, 125}
\definecolor{cautioncolor}{RGB}{195, 85, 70}
\definecolor{rulegray}{RGB}{210, 214, 220}
\definecolor{codebackground}{RGB}{249, 250, 252}

\usepackage{minted}
\usemintedstyle{friendly}
\setminted{
  fontsize=\footnotesize,
  bgcolor=codebackground,
  frame=leftline,
  framerule=0.8pt,
  rulecolor=\color{accentcopper},
  breaklines=true,
  breakanywhere=true,
  tabsize=2,
  linenos=false,
  autogobble=true,
}

\usepackage{tcolorbox}
\tcbuselibrary{skins, breakable}

\newtcolorbox{keypoint}{
  enhanced, breakable, sharp corners, arc=0pt, frame hidden, boxrule=0pt,
  colback=accentcopper!6,
  borderline west={2pt}{0pt}{accentcopper},
  top=6pt, bottom=6pt, left=12pt, right=10pt,
  before skip=1em, after skip=0.9em,
}
\newtcolorbox{caution}{
  enhanced, breakable, sharp corners, arc=0pt, frame hidden, boxrule=0pt,
  colback=cautioncolor!6,
  borderline west={2pt}{0pt}{cautioncolor},
  top=6pt, bottom=6pt, left=12pt, right=10pt,
  before skip=1em, after skip=0.9em,
}

%% ----- 표 -----
\usepackage{booktabs}
\usepackage{tabularx}
\usepackage{array}

%% ----- 섹션 스타일 (간결) -----
\usepackage{titlesec}
\titleformat{\section}[block]
  {\sffamily\bfseries\Large\color{deepblue}}{\thesection}{0.6em}{}
  [\vspace{2pt}{\color{rulegray}\hrule height 0.5pt}\vspace{4pt}]
\titleformat{\subsection}
  {\sffamily\bfseries\normalsize\color{accentcopper}}{\thesubsection}{0.6em}{}

%% ----- 헤더 / 푸터 -----
\usepackage{fancyhdr}
\pagestyle{fancy}
\fancyhf{}
\renewcommand{\headrulewidth}{0pt}
\renewcommand{\footrulewidth}{0pt}
\fancyfoot[L]{\sffamily\small\color{deepblue}\bfseries ARES 코드 변경 이력 · BLE · DC 모터 · BATCH · FinalBuild · 12차시 웹 · 인라인 자산 shim}
\fancyfoot[R]{\sffamily\small\thepage}

\usepackage{hyperref}
\hypersetup{
  colorlinks=true, linkcolor=deepblue, urlcolor=deepblue,
  pdftitle={ARES 코드 수정 이력 — 2026-05-19 이후 누적},
  pdfauthor={(주)코리아사이언스, 동명대학교 게임학부},
}

\renewcommand{\contentsname}{목차}
\renewcommand{\figurename}{그림}
\renewcommand{\tablename}{표}

\title{%
  {\sffamily\bfseries\fontsize{20}{24}\selectfont\color{deepblue}ARES 코드 수정 이력}\\[4pt]
  {\rmfamily\normalsize\color{accentcopper}BLE 명령 파이프라인 · 부팅 안정화 · DC 모터 대칭 · BATCH 블록 · 오프라인 빌드 자동화 · 12차시 × 4미션 웹 재편 · 인라인 자산 shim}\\[2pt]
  {\rmfamily\footnotesize\color{rulegray}\textit{2026-05-19 이후 누적 (롤링 이력)}}
}
\author{\sffamily (주)코리아사이언스\quad/\quad 동명대학교 게임학부}
\date{2026년 5월 19일 \textendash{} 5월 22일}

\begin{document}
\maketitle
\thispagestyle{plain}

\section{개요}

본 보고서는 2026년 5월 19일과 5월 20일에 적용된 ARES 펌웨어 및 웹 코드의
변경 사항을 하나의 합본으로 정리한 자료이다. 두 일자의 변경은 형식적으로는
구분되지만 실질적으로 연결되어 있으며, 5월 19일에 적용한 BLE 명령
파이프라인 최적화의 \emph{재적용 과정}에서 부팅 단계 크래시와 SYS\_SET
핸들러 누락이 드러났고, 같은 흐름에서 DC 모터 출력 비대칭과 Web Bluetooth
호스팅 환경 의존성도 정리되었다.

전체 변경의 큰 줄기는 다음과 같다.

\begin{enumerate}
  \item \textbf{(2026-05-19)} BLE 명령 파이프라인 최적화. ``LED 1번 켜기
        $\rightarrow$ LED 2번 켜기''와 같이 연속된 출력 명령 사이의
        체감 지연(약 250\,ms)을 명령 부류별 응답 정책 분리로 약 40\,ms
        수준까지 단축.
  \item \textbf{(2026-05-20)} 같은 최적화를 재적용하면서 발견된 \emph{펌웨어
        부팅 단계의 OLED 크래시}와 \emph{SYS\_SET 명령 핸들러 누락}을
        교정. 또한 DC 모터의 \texttt{DIR+PWM} 단일 변조 토폴로지에서
        오는 전·후진 출력 비대칭을 듀얼 PWM brake-modulated 구성으로
        근원에서 해소.
  \item \textbf{(2026-05-20)} 학습용 웹앱 \texttt{Web/main.html}이
        \texttt{file://}에서 Web Bluetooth가 동작하지 않는 원인 식별.
        보안 컨텍스트의 문제가 아니라 \emph{ES module의 같은-출처
        정책}이 원인이며, 향후 학생 배포 시 HTTPS 호스팅(또는
        \texttt{localhost})이 사실상 강제됨을 명시.
  \item \textbf{(2026-05-20, 오후)} 명령을 묶어 한 BLE 라운드트립으로
        보내는 BATCH(``한꺼번에 실행'') Blockly 컨테이너 블록 도입과,
        클론 BT05/HM-10 환경에서 멀티 청크 송수신을 안정화하기 위한
        5단계 페이싱/race condition 수정.
  \item \textbf{(2026-05-20, 저녁)} \texttt{file://}에서도 동작하는
        오프라인 배포본을 위해 ES module 8개를 esbuild IIFE 번들로
        합치고 Blockly CDN 사본을 \texttt{Build/vendor/}에 두는
        패키징 절차 정착. 더블클릭 한 번으로 \texttt{Build/} 폴더
        전체를 재생산하는 \texttt{build.bat}/\texttt{build.ps1}
        자동화와 \texttt{Document/FinalBuild.md} 가이드.
  \item \textbf{(2026-05-21)} 학습자 흐름을 1분기 \emph{12차시 × 4미션 =
        48미션} 셸로 전면 재편. \texttt{Web/index.html}(랜딩)
        $\to$ \texttt{Web/main.html}(관제실)의 2단 진입, 차시/미션
        해시 라우터, \texttt{Lesson\{NN\}/lesson.json} + \texttt{overview.html}
        + \texttt{examples/*.xml} 의 런타임 \texttt{fetch()} 도입.
        KoreaScience 데모 페이지(\texttt{ks01.html} 외)는 \texttt{KSWeb/}로
        분리.
  \item \textbf{(2026-05-22)} 21일 재편으로 깨진 오프라인 빌드를
        새 구조에 맞춰 재정착. \texttt{Web/} 소스는 일절 수정하지 않은
        채로, 빌드 시점에 \texttt{overview.html} + 12개 \texttt{lesson.json}
        + \texttt{examples/*.xml}을 한 JS 파일에 인라인하고 \texttt{window.fetch}를
        가로채는 \texttt{vendor/inline\_assets.js} shim을 끼워 \texttt{file://}
        에서도 모든 \texttt{fetch()}가 성공하도록 함. \texttt{build.ps1}을
        5단계 $\to$ 7단계(인라인 shim 생성 단계 추가)로 확장.
\end{enumerate}

\begin{keypoint}
2026-05-19의 작업이 \emph{기능적 개선}이라면, 2026-05-20의 작업은 그
개선이 \emph{재현 가능하게 동작}하도록 만든 안정화이다. 같은 코드라도
하나의 회피 가능한 크래시 때문에 정상이 비정상으로 보이고, 같은 모터라도
PWM 와이어링 한 가지가 사용자 경험을 크게 좌우한다.
\end{keypoint}

\section{2026-05-19 --- BLE 명령 파이프라인 최적화}

ARES 로버의 블록 코딩 사용 중 ``LED 1번 켜기 $\rightarrow$ LED 2번 켜기''와 같이
연속된 명령 사이에 체감 가능한 지연이 발생하던 문제를 분석하여, 명령당
누적 약 250\,ms에 달하던 지연을 출력 명령 기준 약 40\,ms 수준으로 줄였다.
변경 대상은 \texttt{Web/commandexecutor.js}와 \texttt{Pico/main.py}이며,
명령의 의미에 따라 응답 라운드트립을 \emph{선택적으로} 생략하는 정책을
양쪽에 일관되게 도입하였다.

\subsection{문제 진단}

단일 LED 명령 한 번이 종단 간에 소요하는 시간 예산은 다음과 같이
구성되어 있었다.

\begin{table}[h]
\centering
\small
\begin{tabularx}{\linewidth}{@{}llXr@{}}
\toprule
\# & 위치 & 원인 & 시간 \\
\midrule
1 & \texttt{commandexecutor.js:199} & \texttt{sendData(cmd, true)}로 모든 명령에 응답 대기 & 라운드트립 \\
2 & \texttt{bluetooth.js:396} & \texttt{CHUNK\_DELAY=50ms} (\texttt{constants.js:14}) & 50ms \\
3 & \texttt{main.py:86} & \texttt{RECEIVE\_TIMEOUT\_MS=50ms} 무조건 대기 & 50ms \\
4 & \texttt{commandexecutor.js:213} & 모든 명령 후 \texttt{setTimeout(100)} & 100ms \\
5 & BLE GATT  & connection interval (HM-10 기본) & 30--70ms \\
\bottomrule
\end{tabularx}
\caption{단일 LED 명령의 누적 지연(개선 전) --- 합계 약 250\,ms.}
\end{table}

\begin{keypoint}
정상 동작이라도 누적 지연이 \textbf{약 250\,ms / 명령}에 달했다. 카운트
다운, 파티 조명처럼 짧은 간격의 LED 시퀀스가 ``끊기는'' 체감의 주된
원인이다.
\end{keypoint}

\subsection{변경 정책 --- 명령별 분리}

명령을 두 부류로 나누어 응답 대기 여부와 cooldown을 다르게 적용한다.

\begin{table}[h]
\centering
\small
\begin{tabularx}{\linewidth}{@{}lXl@{}}
\toprule
부류 & 명령 & 응답 대기 \\
\midrule
Fire-and-forget & \texttt{LED\_ON}, \texttt{LED\_OFF}, \texttt{MAIN\_LED\_*}, \texttt{[v0 v1 v2 v3 v4]} (LED 패턴),
\texttt{MSG}, \texttt{CLEAR\_DISPLAY}, \texttt{SERVO\_FORWARD}/\texttt{BACKWARD}/\texttt{LEFT}/\texttt{RIGHT}/\texttt{STOP} (연속),
\texttt{DC\_FORWARD}/\texttt{BACKWARD}/\texttt{STOP} (연속), \texttt{GUN\_FIRE} & 아니오 \\
\midrule
응답 대기 유지 & \texttt{BUZZER\_ON,F,D} (D초 blocking), \texttt{SERVO\_t*}/\texttt{DC\_t*} (N초 blocking),
\texttt{SLEEP,N} (N초 blocking), \texttt{DISTANCE}, \texttt{MAGNET}, \texttt{PING} & 예 \\
\bottomrule
\end{tabularx}
\caption{명령 부류와 응답 대기 정책. Pico에서 blocking으로 처리되거나
값을 반환하는 명령은 응답 대기를 유지하여 타이밍과 값 정합성을 보존한다.}
\end{table}

\subsection{Web 측 변경 --- \texttt{Web/commandexecutor.js}}

\subsubsection*{핵심 코드}

\begin{minted}{javascript}
// 즉시 완료 명령 — Pico가 blocking 처리하지 않으므로 응답 대기 불필요.
// 시간 의존적(BUZZER_ON, SERVO_t*, DC_t*, SLEEP) 또는 값 반환(DISTANCE,
// MAGNET, PING)은 이 집합에 넣지 말 것. 응답 대기로 두어야 타이밍/값이
// 보장된다.
FIRE_AND_FORGET_HEADS: new Set([
  'LED_ON', 'LED_OFF',
  'MAIN_LED_ON', 'MAIN_LED_OFF',
  'MSG', 'CLEAR_DISPLAY',
  'SERVO_FORWARD', 'SERVO_BACKWARD', 'SERVO_LEFT', 'SERVO_RIGHT', 'SERVO_STOP',
  'DC_FORWARD', 'DC_BACKWARD', 'DC_STOP',
  'GUN_FIRE',
]),

_isFireAndForget(command) {
  if (command.startsWith('[')) return true;     // LED 패턴 [v0 v1 v2 v3 v4]
  const head = command.split(',')[0];
  return this.FIRE_AND_FORGET_HEADS.has(head);
},

async sendCommand(command) {
  if (!state.isExecuting) {
    Logger.add('[중단] 실행이 중단되었습니다', 'warning');
    return;
  }
  BluetoothManager.updateStatus('명령 실행 중...', STATUS_COLORS.ORANGE);

  const fireAndForget = this._isFireAndForget(command);

  try {
    await BluetoothManager.sendData(command, !fireAndForget);
    if (DEBUG) Logger.add(`[완료] ${command}`, 'info');
  } catch (error) {
    // ... 에러 처리 ...
  }

  // 송신 간격 가드.
  // - fire-and-forget: UART(9600 baud)/BLE 버퍼 오버플로우 방지 최소 마진.
  // - 응답 대기: 이미 라운드트립을 거쳤으므로 가드 단축.
  const cooldown = fireAndForget ? 40 : 20;
  await new Promise(resolve => setTimeout(resolve, cooldown));
}
\end{minted}

\subsubsection*{효과}

\begin{itemize}
  \item 출력 명령(LED, SERVO 연속, DC 연속 등): BLE 라운드트립과 100\,ms
        cooldown이 모두 사라짐 $\rightarrow$ 명령당 약 40\,ms.
  \item 응답 대기 명령: cooldown만 100\,ms$\rightarrow$20\,ms 단축. 의미적
        동기화는 그대로 유지.
\end{itemize}

\subsection{Pico 측 변경 --- \texttt{Pico/main.py}}

웹의 분류와 \emph{동일 기준}으로 Pico에서도 응답 송신 여부를 결정한다.
또한 단일 chunk 명령에 대한 무조건 50\,ms 대기를 제거한다.

\subsubsection*{\texttt{NO\_RESPONSE\_CMDS} 및 \texttt{\_needs\_response}}

\begin{minted}{python}
class AresRover:
    # 응답 송신 생략 명령 (Web/commandexecutor.js의 FIRE_AND_FORGET_HEADS와 동기화).
    # 즉시 완료되는 출력 명령은 응답 라운드트립을 생략하여 처리량을 높이고,
    # 응답 매칭 어긋남(LED_ON 응답이 다음 명령 슬롯에 잘못 들어가는 현상)을 차단한다.
    NO_RESPONSE_CMDS = (
        "LED_ON", "LED_OFF", "MAIN_LED_ON", "MAIN_LED_OFF",
        "MSG", "CLEAR_DISPLAY",
        "SERVO_FORWARD", "SERVO_BACKWARD", "SERVO_LEFT", "SERVO_RIGHT", "SERVO_STOP",
        "DC_FORWARD", "DC_BACKWARD", "DC_STOP",
        "GUN_FIRE",
    )

    def _needs_response(self, data):
        """응답 송신 여부. fire-and-forget 명령은 False (NO_RESPONSE_CMDS 참조)."""
        if data.startswith("["):           # LED 패턴 [v0 v1 v2 v3 v4]
            return False
        head = data.split(",", 1)[0]
        return head not in self.NO_RESPONSE_CMDS
\end{minted}

\subsubsection*{\texttt{\_read\_uart\_line()} 폴링 방식 교체}

기존 구현은 매 호출마다 \texttt{utime.sleep\_ms(RECEIVE\_TIMEOUT\_MS)}로
50\,ms를 무조건 대기한 뒤 두 번째 read를 수행했다. newline이 이미
도착한 단일 chunk 명령에도 적용되어 처리량 천장이 약 16--17 cmd/s에
묶여 있었다. 폴링으로 교체하여 newline 발견 즉시 종료하도록 한다.

\begin{minted}{python}
def _read_uart_line(self):
    """UART에서 완전한 라인 읽기.
    newline 발견 즉시 종료 — 단일 chunk 명령에서 무조건 50ms를 기다리지 않는다.
    멀티 chunk 명령(LED 패턴, SYS_SET 등)은 RECEIVE_TIMEOUT_MS까지 폴링하여 안전.
    """
    if not self.uart.any() and not self.rx_buffer:
        return None

    start = utime.ticks_ms()
    while utime.ticks_diff(utime.ticks_ms(), start) < RECEIVE_TIMEOUT_MS:
        if self.uart.any():
            chunk = self.uart.read()
            if chunk:
                self.rx_buffer += chunk.decode('utf-8', 'ignore')

            # 버퍼 오버플로우 방지
            if len(self.rx_buffer) > MAX_BUFFER_SIZE:
                print("[UART] 버퍼 오버플로우, 초기화")
                self.rx_buffer = ""
                return None

            # newline 즉시 종료
            if '\n' in self.rx_buffer:
                newline_idx = self.rx_buffer.index('\n')
                complete_line = self.rx_buffer[:newline_idx].strip()
                self.rx_buffer = self.rx_buffer[newline_idx + 1:]
                return complete_line
        else:
            utime.sleep_ms(2)

    # 타임아웃: 라인 종료자 없는 부분 데이터 반환
    if self.rx_buffer:
        data = self.rx_buffer.strip()
        if data and not data.endswith(','):
            self.rx_buffer = ""
            return data

    return None
\end{minted}

\subsubsection*{\texttt{run()}의 응답 조건부 송신}

\begin{minted}{python}
# 변경 전
else:
    response = self._process_command(data)
    self._send_response(response)

# 변경 후
else:
    response = self._process_command(data)
    if self._needs_response(data):
        self._send_response(response)
\end{minted}

\subsection{기대 효과 (정량)}

\begin{table}[h]
\centering
\small
\begin{tabular}{@{}lcc@{}}
\toprule
명령 종류 & 변경 전 & 변경 후 \\
\midrule
\texttt{LED\_ON} 등 fire-and-forget & 약 250\,ms & 약 40\,ms \\
\texttt{BUZZER\_ON,F,D} 등 blocking & $D + 250$\,ms & $D + 20$\,ms \\
\texttt{DISTANCE}/\texttt{MAGNET} 등 값 반환 & 약 250\,ms & RTT $+$ 20\,ms \\
Pico 처리량(단일 chunk) & $\approx 16$--$17$ cmd/s & $\approx 100^+$ cmd/s \\
\bottomrule
\end{tabular}
\caption{명령 부류별 종단 간 지연 비교 (단위: ms).}
\end{table}

추가로 응답 매칭 어긋남(이전 fire-and-forget 명령의 응답이 다음
응답 대기 명령의 \texttt{pendingResolve} 슬롯에 잘못 들어가는 코너
케이스) 위험이 \emph{구조적으로} 제거된다. Pico에서 응답 자체를
보내지 않으므로 침범할 응답이 존재하지 않는다.

\begin{caution}
변경된 두 파일은 \textbf{같이 갱신}되어야 한다. 한쪽만 갱신되면:
\begin{itemize}
  \item Web만 갱신: 출력 명령에 응답 대기를 안 하지만 Pico는 여전히
        응답을 보냄 $\rightarrow$ BLE notify 누적, 응답 매칭 어긋남 위험 잔존.
  \item Pico만 갱신: Web이 출력 명령에 응답 대기를 하지만 Pico가 응답을
        안 보냄 $\rightarrow$ 매 명령 5초 timeout.
\end{itemize}
\end{caution}

\section{2026-05-20 --- 재적용과 부팅 안정화}

\subsection{재적용 배경}

전일 적용된 BLE 파이프라인 최적화(\texttt{7e5b93f})는 적용 직후 BLE
연결 자체는 이루어지지만 \texttt{GET\_STATE} 등 모든 명령이 응답하지
않는 현상이 보고되어 \texttt{110bb27}로 revert 되어 있었다.

재현 후 원인 분석한 결과, BLE 파이프라인 최적화 자체는 문제가 없었다.
진짜 원인은 \texttt{Pico/main.py}의 \texttt{boot()}에서 OLED 초기화
실패 시 \texttt{robot.oled = None} 상태인데 \texttt{None.fill(0)}이
호출되어 \texttt{AttributeError}가 발생하고, 이 예외가 메인 루프 진입을
원천적으로 막아버린 데 있었다. 즉 \emph{펌웨어 부팅 단계의 무관한
크래시}가 BLE 파이프라인 변경의 효과를 가린 것이다.

\subsection{관찰의 단서 --- BLE 연결은 되는데 응답이 없다}

BLE 모듈(BT05/HM-10)은 Pico의 UART에 연결된 외부 모듈로, BLE
페어링/연결은 모듈 자체 펌웨어가 처리하고 Pico는 UART로 데이터만
주고받는다. 따라서 Pico의 \texttt{main.py}가 boot에서 죽어도 BLE 링크는
멀쩡하다. 외부에서 보이는 증상은 ``연결은 되는데 어떤 명령도
응답하지 않는다''로, 이는 정확히 \emph{Pico가 메인 루프에 진입조차
못한} 상태의 결과이다.

\subsection{수정 --- OLED None 가드}

\begin{minted}{python}
# 변경 전 (Pico/main.py boot())
self.robot.oled.fill(0)
self.robot.oled.text("ARES READY", 0, 0)
self.robot.oled.show()

# 변경 후
if self.robot.oled is not None:
    self.robot.oled.fill(0)
    self.robot.oled.text("ARES READY", 0, 0)
    self.robot.oled.show()
\end{minted}

가드 추가 후 OLED 모듈이 비활성(또는 I2C 초기화 실패) 상태여도
\texttt{boot()}가 완료되며, 메인 루프가 정상 진입한다. 이 가드는 다른
하드웨어 모듈에 적용된 \texttt{if robot.module:} 패턴과 일관된
방어 코드이다.

가드 적용 후 5월 19일의 BLE 최적화를 동일 형태로 재적용하였고
(\texttt{600243a}), 정상 동작을 확인하였다.

\section{2026-05-20 --- SYS\_SET 핸들러 누락 수정}

\subsection{증상}

대시보드에서 ``최대 속도'' 슬라이더를 낮추고 ``Pico에 업로드'' 버튼을
누르면 ``설정이 Pico에 업로드되었습니다!'' 라는 알림이 떴지만, 다른
탭으로 갔다가 돌아오면 슬라이더가 다시 80\%로 복귀했다.

\subsection{원인}

\texttt{Pico/process\_data.py}의 명령 디스패처에서:

\begin{minted}{python}
elif data.startswith("SYS_SET"):
    return self._handle_sys_set(data)
\end{minted}

호출 대상인 \texttt{\_handle\_sys\_set} 메서드가 \emph{정의되어 있지
않았다}. 메서드 헤더(\texttt{def \_handle\_sys\_set(self, data):})가
빠져 있어, 핸들러 본문이 직전의 \texttt{\_handle\_set\_module}의
\texttt{return} 뒤에 떨어져 \emph{도달 불가능한 코드(dead code)} 가
되어 있었다. SYS\_SET 명령이 도달하면 \texttt{AttributeError}가 발생
하여 \texttt{main.py}의 예외 처리에 잡혀 \texttt{"ERROR"}가 반환되고,
\texttt{system.json}은 갱신되지 않았다. Web 측은 응답을 확인하지 않고
즉시 alert를 띄우므로 사용자에게는 ``저장된 것처럼'' 보였다.

\subsection{수정}

\begin{minted}{python}
def _handle_sys_set(self, data):
    """시스템 설정 저장: SYS_SET,max_speed,col_dist,auto_stop,name"""
    try:
        parts = data.split(',')
        if len(parts) >= 5:
            max_speed = parts[1]
            col_dist  = parts[2]
            auto_stop = parts[3]
            name      = ",".join(parts[4:])
            sys_config.save_config(max_speed, col_dist, auto_stop, name)
            return 1
    except Exception as e:
        print(f"SYS_SET 오류: {e}")
    return 0
\end{minted}

수정 후 슬라이더 값이 \texttt{system.json}에 영구 저장되고, Pico 재부팅
후에도 유지된다.

\begin{caution}
Web 측 \texttt{saveSystemSettings()}는 SYS\_SET 응답을 검증하지 않고
무조건 성공 alert를 띄운다. 후속 개선 항목으로 \texttt{sendData(cmd,
true)}로 응답 대기 후 결과를 분기하는 작업을 분리하여 진행한다.
\end{caution}

\section{2026-05-20 --- DC 모터 H-브릿지 비대칭 진단과 듀얼 PWM 대칭화}

\subsection{증상의 변천}

\begin{enumerate}
  \item 초기: 전진이 후진보다 명백히 약함(특히 \texttt{max\_speed}가
        높을수록 격차가 큼).
  \item 단순 PWM 반전 적용 후: 80\,\%에서는 양방향 정상이나, 슬라이더
        55\,\%\,이하에서 후진이 시동되지 않음.
  \item 후진에 시동 데드존 lift 적용(56\,\%): 시동은 되지만 후진 토크가
        56--60\,\% 구간에서 약함.
  \item 모터 리드를 물리적으로 뒤집는 진단 실험: 회전 방향만 반대로
        바뀌고 전·후진 속도 특성은 함수 호출 기준 그대로 유지됨.
\end{enumerate}

\subsection{진단 결과 --- 비대칭은 모터가 아니라 드라이버 입력 토폴로지}

기존 코드는 \texttt{DIR + PWM} 한 쌍으로 H-브릿지를 구동하였다. 이때
드라이버 칩의 \texttt{IN1}/\texttt{IN2} 두 입력 중 한쪽은 디지털 DIR로
고정되고 다른 쪽만 PWM이 들어가므로, 두 회전 방향이 \emph{서로 다른
PWM 변조 모드}로 동작한다.

\begin{table}[h]
\centering
\small
\begin{tabularx}{\linewidth}{@{}lXX@{}}
\toprule
모드 & 동작 & 특성 \\
\midrule
\texttt{DIR=1}, IN2 PWM (전진) & drive$\leftrightarrow$\textbf{brake} 변조 & 권선 전류 연속, 시동 우수, 듀티가 작을수록 토크 큼 \\
\texttt{DIR=0}, IN2 PWM (후진) & drive$\leftrightarrow$\textbf{coast} 변조 & 권선 전류 불연속, 시동 데드존 존재, 약 56\,\% 이하에서 시동 불가 \\
\bottomrule
\end{tabularx}
\caption{단일 \texttt{DIR+PWM} 토폴로지에서 회전 방향별 PWM 변조 모드.}
\end{table}

모터 리드를 뒤집어도 함수 단위의 비대칭이 그대로 유지된 실험 결과가
이 가설을 결정적으로 확증한다. 비대칭은 \emph{드라이버 입력 측 패턴}에
귀속되며, 모터 권선이나 기어박스의 마찰 비대칭은 아니다.

\subsection{듀얼 PWM brake-modulated 구동}

RP2040에서 \texttt{GP8}/\texttt{GP9}는 같은 PWM 슬라이스(PWM4)의 A/B
채널이다. 두 핀을 모두 PWM으로 설정한 뒤, 한쪽을 \texttt{PWM\_MAX}로
고정하고 반대쪽을 \texttt{PWM\_MAX - speed}로 변조하면 양방향 모두
brake-modulated가 되어 토크 곡선이 대칭화된다.

\begin{minted}{python}
# Pico/dcmotor.py
from machine import Pin, PWM
from pins import DCMOTOR_DIR_PIN, DCMOTOR_PWM_PIN

IN1_PIN = DCMOTOR_DIR_PIN  # GP8 (PWM4A)  — 이름은 레거시 호환 유지
IN2_PIN = DCMOTOR_PWM_PIN  # GP9 (PWM4B)

PWM_FREQUENCY = 20000
PWM_STOP = 0
PWM_MAX  = 65535


class KSDCMotor:
    """DC 모터 컨트롤러 (양방향 brake-modulated H-브릿지)."""

    def __init__(self):
        self.IN1 = PWM(Pin(IN1_PIN))
        self.IN1.freq(PWM_FREQUENCY)
        self.IN1.duty_u16(PWM_STOP)

        self.IN2 = PWM(Pin(IN2_PIN))
        self.IN2.freq(PWM_FREQUENCY)
        self.IN2.duty_u16(PWM_STOP)

        self._is_running = False

    def dc_forward(self, speed):
        """전진: IN1 고정 HIGH, IN2를 (PWM_MAX-speed)로 변조."""
        if speed <= 0:
            self.dc_stop(); return
        self.IN1.duty_u16(PWM_MAX)
        self.IN2.duty_u16(PWM_MAX - speed)
        self._is_running = True

    def dc_backward(self, speed):
        """후진: IN2 고정 HIGH, IN1을 (PWM_MAX-speed)로 변조."""
        if speed <= 0:
            self.dc_stop(); return
        self.IN1.duty_u16(PWM_MAX - speed)
        self.IN2.duty_u16(PWM_MAX)
        self._is_running = True

    def dc_stop(self):
        """정지 (coast: 둘 다 LOW)."""
        self.IN1.duty_u16(PWM_STOP)
        self.IN2.duty_u16(PWM_STOP)
        self._is_running = False
\end{minted}

\subsection{기대 효과}

\begin{table}[h]
\centering
\small
\begin{tabularx}{\linewidth}{@{}lXXX@{}}
\toprule
슬라이더 & 전진 실효 듀티 & 후진 실효 듀티 & 비고 \\
\midrule
100\,\% & 100\,\% & 100\,\% & 풀파워, 양방향 동일 \\
80\,\%  & 80\,\%  & 80\,\%  & 동일 \\
50\,\%  & 50\,\%  & 50\,\%  & 동일 \\
20\,\%  & 20\,\%  & 20\,\%  & 동일 (시동은 모터 마찰 한계에 의존) \\
0\,\%   & 0       & 0       & coast \\
\bottomrule
\end{tabularx}
\caption{듀얼 PWM brake-modulated 구동의 양방향 대칭 출력.}
\end{table}

\begin{keypoint}
시동 데드존(\texttt{BACKWARD\_DEADZONE}) lift 같은 보정 hack이 더는
필요하지 않다. \emph{소프트웨어로 우회}하던 비대칭을 \emph{펌웨어
구성 변경}으로 \emph{근원에서} 제거한 사례이다. 회로 배선이나 드라이버
교체 없이 펌웨어만으로 해결되었다는 점이 특히 가치 있다.
\end{keypoint}

\subsection{잔여 사항}

\begin{itemize}
  \item \texttt{pins.py} 및 \texttt{system\_config.py}의 핀 이름
        (\texttt{DCMOTOR\_DIR\_PIN}, \texttt{pin\_dcmotor\_dir})은
        의미상 ``DIR''이 아니지만 외부 API(\texttt{SET\_PIN} 명령
        등)와의 호환을 위해 레거시 명칭을 유지했다. 향후 정리 시
        \texttt{IN1}/\texttt{IN2} 또는 \texttt{A}/\texttt{B}로
        명명 통일이 권장된다.
  \item 모터의 정적 마찰 한계로 약 20\,\% 이하의 매우 낮은 듀티에서는
        여전히 회전하지 않을 수 있다. 이는 회로/펌웨어 문제가 아닌
        모터 물성으로, 학습 환경에서는 \texttt{max\_speed}를 30\,\%
        이상에서 운영할 것을 권장한다.
\end{itemize}

\section{2026-05-20 --- Web Bluetooth와 HTTPS 서버 사용 사유}

\subsection{관찰}

\texttt{Web/ks01.html}은 \texttt{file://} 프로토콜로 직접 열어도 BLE
연결이 가능했다. 반면 정식 학습용 \texttt{Web/main.html}은
\texttt{file://}에서 연결되지 않고, \texttt{http://localhost} 또는
\texttt{https://} 호스팅 환경에서만 동작한다.

\subsection{원인 --- Web Bluetooth가 아니라 ES module의 same-origin 정책}

직관적으로는 ``Web Bluetooth가 \texttt{file://}에서 secure context로
인정되지 않기 때문''으로 생각하기 쉽지만, 실제 원인은 다르다. 두
파일의 스크립트 로딩 방식이 결정적으로 다르다.

\begin{table}[h]
\centering
\small
\begin{tabularx}{\linewidth}{@{}lXX@{}}
\toprule
파일 & 스크립트 로딩 방식 & \texttt{file://}에서 동작? \\
\midrule
\texttt{ks01.html} & 모든 코드가 인라인 \texttt{<script>...</script>} & 정상 동작 \\
\texttt{main.html} & \texttt{<script type="module" src="main.js">}, \texttt{main.js}는 \texttt{bluetooth.js} 등 다수 모듈을 \texttt{import} & 차단 (모듈 fetch 실패) \\
\bottomrule
\end{tabularx}
\caption{스크립트 로딩 방식과 \texttt{file://} 동작 가부.}
\end{table}

Chromium 기반 브라우저는 \texttt{file://}을 secure context로 인정하므로
\texttt{navigator.bluetooth} API 자체는 \texttt{file://}에서도 호출
가능하다. 그러나 \texttt{<script type="module">}은 \emph{같은 출처
(same-origin) 정책}을 따르는데, 모든 \texttt{file://} URL은 같은
디렉터리에 있더라도 \texttt{origin: null}로 취급되어 서로 다른 출처로
간주된다. 따라서 \texttt{main.js}를 비롯한 ES 모듈 fetch가 CORS
정책으로 차단되고, 콘솔에는 다음과 같은 오류가 기록된다.

\begin{minted}{text}
Access to script at 'file:///.../main.js' from origin 'null'
has been blocked by CORS policy: Cross origin requests are
only supported for protocol schemes: http, data, isolated-app,
chrome-extension, chrome-untrusted, https, chrome.
\end{minted}

ES 모듈이 로드되지 않으면 \texttt{BluetoothManager} 코드는 실행될
기회조차 없다. 즉 사용자가 본 ``BLE가 연결되지 않는'' 증상은 Web
Bluetooth API의 보안 제한이 아니라 \emph{BLE 코드가 아예 실행되지
않은} 결과이다. 증상이 같아 보였을 뿐이다.

\begin{caution}
\texttt{ks01.html}이 \texttt{file://}에서 동작하는 것은 데스크톱 Chrome
환경에 한정될 가능성이 크다. Android Chrome은 \texttt{file://}에서 Web
Bluetooth를 더 엄격히 제한하는 경우가 있으며, 모바일 배포를 가정한다면
이 경로에 의존해서는 안 된다.
\end{caution}

\subsection{해결 옵션과 권장}

\begin{table}[h]
\centering
\small
\begin{tabularx}{\linewidth}{@{}lXl@{}}
\toprule
방법 & 설명 & 권장 \\
\midrule
A. \texttt{python -m http.server} & \texttt{localhost}는 secure context로 인정되고 ES 모듈도 같은 출처에서 fetch 가능 & 개발 단계 \\
B. HTTPS 호스팅 (GitHub Pages 등) & 어디서나 동작, 모바일 포함. \texttt{file://} 의존 사라짐 & 학생 배포 \\
C. 단일 HTML로 번들링(esbuild/rollup) & \texttt{main.js}와 모듈을 인라인 IIFE로 합쳐 \texttt{ks01.html}과 같은 형태 & 오프라인 단일 파일이 꼭 필요한 경우 \\
D. ES module 제거, 전역 객체로 통합 & \texttt{import/export} 제거 후 인라인 \texttt{<script>}로 합치기 & 권장하지 않음 (모듈성 손실 큼) \\
\bottomrule
\end{tabularx}
\caption{\texttt{main.html}의 \texttt{file://} 미동작에 대한 해결 옵션.}
\end{table}

\begin{keypoint}
학생 배포를 전제하면 \textbf{B(HTTPS 호스팅)}\,가 가장 깔끔하다. 모바일
브라우저 호환성, 즐겨찾기/PWA 설치, 자동 업데이트 측면 모두 유리하다.
\texttt{ks01.html}처럼 단일 파일 배포가 정책적으로 요구되는 경우에는
\textbf{C(번들링)}\,가 정공법이다. 어느 쪽이든 \texttt{file://} 직접
열기에 의존하는 워크플로우는 장기적으로 권장되지 않는다.
\end{keypoint}

\section{2026-05-20 (오후) --- BATCH ``한꺼번에 실행'' 명령 블록과 파이프라인 정착}

BLE 파이프라인을 출력 명령에 대해 fire-and-forget으로 최적화한 결과
단일 명령당 약 40\,ms의 Web 측 cooldown까지는 줄였지만, 명령이 수십 개에
달하는 시퀀스(카운트다운, 사이렌, LED 시퀀스)에서는 누적 cooldown이
여전히 수 초에 이르는 문제가 남았다. 본 절은 명령을 묶어 한 번의 BLE
라운드트립으로 보내는 BATCH 메커니즘과, 이를 클론 BT05/HM-10 환경에서
안정적으로 동작시키기 위해 진행한 단계별 수정을 정리한다.

\subsection{동기와 설계}

명령을 묶을 컨테이너 블록을 Blockly에 도입하고, 안에 들어간 자식 블록들의
명령 문자열을 \texttt{|}로 이어 \texttt{BATCH;cmd1|cmd2|cmd3} 형태의 단일
명령으로 Pico에 보낸다. Pico는 자체 큐로 순차 실행하고 마지막에 단일
응답을 돌려준다.

\begin{table}[h]
\centering
\small
\begin{tabularx}{\linewidth}{@{}lXX@{}}
\toprule
이득 & 의존 매개변수 & BATCH로 달성 \\
\midrule
명령당 cooldown 누적 제거 & Web cooldown 정책 & 단일 cooldown(20\,ms)으로 압축 \\
\texttt{SLEEP}/\texttt{BUZZER\_ON} 타이밍 정확도 & Pico 내부 처리 & 네트워크 지연 영향 제거 \\
응답 매칭 어긋남 위험 제거 & 단일 ACK 구조 & sub-command별 ACK 없음 \\
BLE 송신 시간 단축 & 청크 수 & 부수적 (\texttt{CHUNK\_DELAY}에 영향) \\
\bottomrule
\end{tabularx}
\caption{BATCH 도입의 본질적 이득 분해.}
\end{table}

\subsection{Blockly 블록 디자인과 차단 정책}

새 블록 \texttt{batch\_block}은 단일 \texttt{statement input}을 가지는
컨테이너이며 보라색(\#8E44AD)으로 시각적으로 두드러진다. 안에 들어가면
안 되는 블록(센서 값 반환, 반복/조건/변수, 함수 호출, 자기 자신)은
미리 집합으로 명시한다.

\begin{minted}{javascript}
// Web/blocklyconfig.js
export const BATCH_FORBIDDEN_TYPES = new Set([
  // 값 반환
  'check_distance', 'check_magnetic', 'pico_check_device',
  // 제어 흐름
  'controls_if', 'controls_whileUntil', 'controls_repeat_ext',
  // 변수/함수 (제어 흐름은 Web 측 책임)
  'variables_set', 'assign_variable', 'math_change',
  'procedures_callnoreturn', 'procedures_callreturn',
  'procedures_defnoreturn', 'procedures_defreturn',
  // 중첩 금지
  'batch_block',
]);
\end{minted}

학생이 금지 블록을 batch 안에 떨어뜨리면 \texttt{setOnChange} 콜백이
자식을 순회해 노란 ⚠️ 경고를 즉시 띄운다. 실행 시점에도
\texttt{commandexecutor.js}의 \texttt{\_processBatch}가 한 번 더 검증하여
금지 블록 발견 시 빨간 로그와 함께 실행을 중단한다(이중 방어).

\subsection{프로토콜과 핵심 구현}

\begin{minted}{javascript}
// Web/commandexecutor.js
async _processBatch(block) {
  const commands = [];
  let cur = block.getInputTargetBlock('DO');
  while (cur) {
    if (BATCH_FORBIDDEN_TYPES.has(cur.type)) {
      Logger.add(`[오류] '${cur.type}' 블록은 [한꺼번에 실행] 안에 넣을 수 없습니다.`, 'error');
      state.isExecuting = false;
      return;
    }
    const cmd = this.generateCommand(cur);
    if (cmd) commands.push(cmd);
    cur = cur.getNextBlock();
  }
  if (commands.length === 0) return;
  const payload = `BATCH;${commands.join('|')}`;
  await BluetoothManager.sendData(payload, true);   // 응답 대기 1회
  await new Promise(r => setTimeout(r, 20));
}
\end{minted}

\begin{minted}{python}
# Pico/process_data.py
def _handle_batch(self, data):
    """BATCH;cmd1|cmd2|cmd3 — 기존 디스패처로 차례 실행 후 단일 응답."""
    try:
        body = data[len("BATCH;"):]
        for cmd in body.split('|'):
            cmd = cmd.strip()
            if not cmd:
                continue
            self.process(cmd)
        return 1
    except Exception as e:
        print(f"BATCH 오류: {e}")
        return 0
\end{minted}

\subsection{진행 중 발견된 안정화 이슈와 단계별 수정}

초기 구현은 단일 청크에 들어가는 매우 짧은 BATCH에서만 정상 동작했고,
일반 학습 시나리오(LED 2개 + SLEEP 등)는 멀티 청크로 전송되며 다음 5단계의
문제를 차례로 드러냈다.

\subsubsection*{1단계 --- 부분 라인 반환으로 인한 BATCH 절단}

\texttt{\_read\_uart\_line()}이 \texttt{RECEIVE\_TIMEOUT\_MS} 동안 newline을
못 만나면 \emph{쉼표로 끝나지 않은 부분 데이터를 반환}하는 옛 동작 때문에,
첫 청크 \texttt{"BATCH;LED\_ON,0,1|LED"}가 그대로 디스패처에 넘어가
\texttt{LED\_ON,0,1}만 실행되고 나머지가 잘려나갔다. 부분 반환을 제거하고
buffer를 유지한 채 \texttt{None}을 반환하도록 변경했다.

\begin{minted}{python}
# Pico/main.py
def _read_uart_line(self):
    if not self.uart.any() and not self.rx_buffer:
        return None
    start = utime.ticks_ms()
    while utime.ticks_diff(utime.ticks_ms(), start) < RECEIVE_TIMEOUT_MS:
        if self.uart.any():
            chunk = self.uart.read()
            if chunk:
                self.rx_buffer += chunk.decode('utf-8', 'ignore')
            if len(self.rx_buffer) > MAX_BUFFER_SIZE:
                print("[UART] 버퍼 오버플로우, 초기화")
                self.rx_buffer = ""
                return None
            if '\n' in self.rx_buffer:
                idx = self.rx_buffer.index('\n')
                line = self.rx_buffer[:idx].strip()
                self.rx_buffer = self.rx_buffer[idx + 1:]
                return line
        else:
            utime.sleep_ms(2)
    # 타임아웃 시 buffer 유지, 다음 호출에서 이어받음
    return None
\end{minted}

\subsubsection*{2단계 --- 청크 손실 진단과 페이싱 조정}

원인 추적 끝에 BT05/HM-10 클론에서 청크 사이 간격이 BLE connection
interval(30--70\,ms)과 겹쳐 청크가 일부 손실됨을 확인했다. AT 명령으로
connection interval을 단축하는 정식 경로는 클론에서 미지원이므로 Web 측
페이싱을 보수적으로 조정했다.

\begin{table}[h]
\centering
\small
\begin{tabular}{@{}lll@{}}
\toprule
매개변수 & 변경 전 & 정착값 \\
\midrule
\texttt{Web/constants.js : CHUNK\_DELAY} & 50\,ms & \textbf{100\,ms} \\
\texttt{Pico/main.py : RECEIVE\_TIMEOUT\_MS} & 50\,ms & \textbf{500\,ms} \\
\bottomrule
\end{tabular}
\caption{클론 BLE 모듈을 위한 페이싱 정착값.}
\end{table}

\subsubsection*{3단계 --- GATT ACK 왕복 제거}

\texttt{writeValueWithResponse}는 매 청크마다 GATT 수준 ACK를 기다리므로
멀티 청크 송신 시간이 누적된다. characteristic이 지원하면
\texttt{writeValueWithoutResponse}를 우선 사용하고, 그렇지 않으면 기존
방식으로 fallback하도록 했다.

\begin{minted}{javascript}
// Web/bluetooth.js
const useWithoutResponse =
  state.characteristic.properties &&
  state.characteristic.properties.writeWithoutResponse;
// for-loop 내부에서:
if (useWithoutResponse) {
    await state.characteristic.writeValueWithoutResponse(chunk);
} else {
    await state.characteristic.writeValueWithResponse(chunk);
}
\end{minted}

\subsubsection*{4단계 --- 응답 race condition 수정}

BATCH가 Pico에서 매우 빠르게 처리되어 마지막 청크 송신 완료 직후 응답이
도착하는데, \texttt{pendingResolve}가 송신 \emph{후}에 설치되어 있어 응답이
\texttt{state.pendingResolve === null} 상태에서 버려졌다. 5초
\texttt{RESPONSE\_TIMEOUT} 후 명령이 실패한 것처럼 보였지만 Pico에서는
이미 두 LED 모두 점등되어 있었다. promise 설치를 송신 시작 \emph{전}으로
옮겨 race를 닫았다.

\begin{minted}{javascript}
// Web/bluetooth.js — sendData 발췌
let responsePromise = null;
if (waitForResponse) {
    responsePromise = new Promise((resolve, reject) => {
        state.pendingResolve = resolve;     // ← 송신 전 설치
        state.pendingReject = reject;
        state.pendingTimeout = setTimeout(() => {
            state.pendingResolve = null;
            state.pendingReject = null;
            state.pendingTimeout = null;
            reject(new Error('응답 시간 초과'));
        }, BLUETOOTH_CONFIG.RESPONSE_TIMEOUT);
    });
}
// ... 청크 송신 ...
return waitForResponse ? responsePromise : 'OK';
\end{minted}

\subsubsection*{5단계 --- 단일 청크 명령의 마지막 \texttt{CHUNK\_DELAY} 제거}

\texttt{CHUNK\_DELAY=100\,ms}로 정착한 뒤 일반 LED 명령처럼 단일 청크
명령조차 매번 100\,ms를 추가로 기다리는 부작용이 관찰되었다. 송신 루프가
\emph{마지막 청크 송신 후에도} 무조건 \texttt{CHUNK\_DELAY}를 실행했기
때문이다. 청크 \emph{사이}에만 페이싱이 적용되도록 guard를 추가했다.

\begin{minted}{javascript}
for (let i = 0; i < encodedData.length; i += BLUETOOTH_CONFIG.MAX_CHUNK_SIZE) {
    const chunk = encodedData.slice(i, ...);
    if (useWithoutResponse) await state.characteristic.writeValueWithoutResponse(chunk);
    else                    await state.characteristic.writeValueWithResponse(chunk);
    // 청크 사이에만 페이싱 — 마지막 청크 후엔 즉시 종료
    if (i + BLUETOOTH_CONFIG.MAX_CHUNK_SIZE < encodedData.length) {
        await this.delay(BLUETOOTH_CONFIG.CHUNK_DELAY);
    }
}
\end{minted}

\subsection{최종 효과}

\begin{table}[h]
\centering
\small
\begin{tabularx}{\linewidth}{@{}lXXX@{}}
\toprule
명령 형태 & 변경 전 & BATCH 도입 직후 & BATCH 정착 후 \\
\midrule
단일 청크 출력 (\texttt{LED\_ON} 등) & 약 40\,ms & 약 140\,ms (마지막 \texttt{CHUNK\_DELAY} 100\,ms 가산) & 약 40\,ms (마지막 \texttt{CHUNK\_DELAY} guard로 회복) \\
멀티 청크 시퀀스 100명령 누적 cooldown & 약 4\,초 & --- & 단일 cooldown으로 압축 \\
\texttt{SLEEP}/\texttt{BUZZER\_ON} 타이밍 정확도 & 네트워크 지연 영향 & --- & Pico 내부 처리, 정확 \\
응답 매칭 어긋남 & 가능성 있음 & 가능성 있음 & 단일 응답이라 구조적 제거 \\
\bottomrule
\end{tabularx}
\caption{BATCH 도입 전후의 명령별 종단 효과.}
\end{table}

\begin{keypoint}
이번 작업은 두 가지 가치가 함께 간다. 첫째, 카운트다운/사이렌 같은 시간
의존 시퀀스가 \textbf{시각적으로 부드러워졌다}. 둘째, 학생이 ``하나씩 결과를
확인하고 싶은 동작''과 ``빠르게 한 번에 보내고 싶은 동작''을 \emph{Blockly
블록 위에 명시적으로 표현}하게 되었다. 추상적으로는 원격 호출(RPC)과 원격
스크립트 송신의 차이를 초등 수준에서 체험하는 학습 도구이기도 하다.
\end{keypoint}

\subsection{BATCH 관련 잔여 사항}

\begin{itemize}
  \item \texttt{RESPONSE\_TIMEOUT=5\,초}는 BATCH 안의 누적 \texttt{SLEEP}이
        5\,초를 넘으면 timeout이 먼저 도착한다. 후속 개선으로 BATCH의 명령
        합에 비례한 \emph{동적 timeout} 산정 또는 BATCH 전용 더 긴 timeout
        도입을 검토한다.
  \item BATCH 안 명령을 단축 토큰(예: \texttt{LED\_ON,0,1}\,$\rightarrow$\,\texttt{L+0})
        으로 packing하면 한 청크에 더 많은 명령이 들어가 송신 시간을 더
        단축할 수 있다. 디스패처 호환성 작업이 필요하므로 별도 분리.
\end{itemize}

\section{잔여 위험 및 후속 작업}

\subsection{Web 측 \texttt{saveSystemSettings()} 응답 검증}

SYS\_SET 송신 후 무조건 성공 alert를 띄우는 현행 동작은 펌웨어에서
실제로 저장 실패한 경우를 사용자에게 가린다. 응답 대기 명령으로
승격하여 \texttt{"1"}/\texttt{"0"} 또는 \texttt{"OK"}/\texttt{"ERROR"}
구분 후 분기 alert를 띄우는 작업이 권장된다.

\subsection{UART baud rate}

가장 본질적인 처리량 천장은 BLE 모듈 $\leftrightarrow$ Pico 사이의 9600
baud UART이다. HM-10/BT05의 AT 명령(\texttt{AT+BAUD8} 등)으로 38400
또는 115200 baud로 상향하고 \texttt{Pico/main.py}의
\texttt{UART\_BAUDRATE}를 동기화하면 byte당 전송 시간을 4--12배 단축할
수 있다. BLE 모듈 펌웨어 설정과 Pico 코드 동시 수정이 필요하므로 별도
작업으로 분리한다.

\subsection{호스팅 환경 정비}

\texttt{Web/main.html}의 ES module 의존을 그대로 두면서 학생 배포를
원활하게 하려면 GitHub Pages 등 HTTPS 호스팅을 정식 채택할 필요가 있다.
도메인/배포 자동화/PWA manifest 추가가 후속 작업 후보이다.

\subsection{권장 회귀 테스트}

\begin{enumerate}
  \item \textbf{카운트다운} --- \texttt{LED\_OFF} 5개 연속 소등. 명령
        사이 체감 지연이 거의 사라져야 한다.
  \item \textbf{LED 패턴 + 거리 센서 혼합} --- \texttt{[v0 v1 v2 v3 v4]}
        직후 \texttt{DISTANCE}. 거리값이 \texttt{1} 같은 LED 응답으로
        오염되지 않는지 확인.
  \item \textbf{부저 사이렌} --- \texttt{BUZZER\_ON,400,0.5} 와
        \texttt{BUZZER\_ON,800,0.5}의 반복. 타이밍 정확성이 유지되어야
        한다.
  \item \textbf{시간 제한 모터} --- \texttt{DC\_tFORWARD,N},
        \texttt{DC\_tBACKWARD,N}의 N초 동기화가 유지되는지 확인.
  \item \textbf{DC 모터 양방향 대칭} --- \texttt{max\_speed}를 30, 50,
        80, 100\,\%로 바꿔가며 전진/후진 속도가 거의 같게 느껴지는지
        확인.
  \item \textbf{대시보드 \texttt{max\_speed} 영구 저장} --- 슬라이더
        조정 $\rightarrow$ 업로드 $\rightarrow$ 다른 탭 이동 $\rightarrow$ 대시보드 복귀 시
        값이 유지되는지 확인. Pico 재부팅 후에도 유지되어야 한다.
  \item \textbf{BATCH 카운트다운} --- ``한꺼번에 실행'' 블록 안에 LED 5
        $\rightarrow$ 1 점등/소등 + \texttt{SLEEP,1}을 넣고 실행. 한 번의
        \texttt{[전송] BATCH;...} 로그와 단일 \texttt{[묶음 완료]} 응답이
        떠야 하며, LED 시퀀스가 끊김 없이 매끄러워야 한다.
  \item \textbf{BATCH 금지 블록 차단} --- 거리 센서/반복/조건 블록을
        ``한꺼번에 실행'' 안에 끌어놓을 때 노란 ⚠️ 경고가 즉시 표시되고,
        실행 시 빨간 오류 로그와 함께 중단되는지 확인.
  \item \textbf{단일 청크 명령 속도 회복} --- 일반 \texttt{LED\_ON} 단일
        블록을 5\,$\sim$\,10회 연속 실행. 명령 사이 체감 지연이 \texttt{CHUNK\_DELAY}
        도입 전 수준(약 40\,ms)으로 돌아왔는지 확인.
  \item \textbf{오프라인 빌드 재현} --- \texttt{Build/} 폴더 삭제 후
        \texttt{build.bat} 더블클릭. 5단계가 모두 성공하고 \texttt{Build/}가
        새로 생성된 뒤, \texttt{Build/index.html}을 \texttt{file://}로 열어
        한글 표시·BLE 페어링·\texttt{한꺼번에 실행} 시퀀스가 모두 정상
        동작하는지 확인.
\end{enumerate}

\section{2026-05-20 (저녁) --- FinalBuild: 오프라인 (\texttt{file://}) 배포 빌드와 자동화}

학생 PC 배포와 인터넷이 없는 학교 환경을 가정하면 \texttt{Web/main.html}의
ES module + Blockly CDN 의존 구조를 그대로 둘 수 없다. 본 절은
\texttt{main.html}을 \texttt{file://}로 직접 열어도 BLE까지 전 기능이
동작하도록 패키징하는 절차를 정착시킨 작업, 그 과정에서 마주친 함정들,
그리고 \texttt{build.bat} 한 번의 더블클릭으로 재현 가능한 자동화
파이프라인을 정리한다.

\subsection{두 가지 외부 의존성과 해법}

\texttt{file://}에서 \texttt{main.html}이 실패하는 원인은 두 가지였다.

\begin{table}[h]
\centering
\small
\begin{tabularx}{\linewidth}{@{}lXX@{}}
\toprule
의존성 & \texttt{file://}에서의 문제 & 본 빌드의 해법 \\
\midrule
\texttt{Web/} 8개 ES module & \texttt{<script type="module">}이 fetch될 때 \texttt{origin: null}의 같은-출처 정책으로 차단 & esbuild로 단일 IIFE 번들(\texttt{main.bundle.js})로 합쳐 일반 \texttt{<script>}로 로드 \\
Blockly CDN 3개 스크립트 & \texttt{unpkg.com}에 인터넷 연결 필요 & \texttt{Build/vendor/}에 사본 두고 상대 경로로 참조 \\
\bottomrule
\end{tabularx}
\caption{\texttt{file://} 동작을 막던 두 외부 의존성과 본 빌드의 대응.}
\end{table}

\subsection{산출물 --- \texttt{Build/} 폴더 한 단위}

빌드 결과는 프로젝트 루트의 \texttt{Build/} 폴더 한 곳에 모인다. 이
폴더 자체가 \emph{완전한 자족(self-contained) 배포 단위}이며, USB로
복사하거나 zip으로 묶어 옮기면 그 자체로 ARES 앱이 동작한다.

\begin{minted}{text}
Build/                          1.2 MB
├── index.html                  9.9 KB   (main.html 패치: vendor/ 경로 + main.bundle.js defer)
├── main.bundle.js             77.6 KB   (8개 ES module을 합친 IIFE 번들)
├── styles.css                 11.3 KB
├── dashboard.html             32.6 KB   (인라인 스크립트만 사용, 수정 없음)
└── vendor/
    ├── blockly_compressed.js  900.8 KB
    ├── blocks_compressed.js    86.8 KB
    └── python_compressed.js    50.3 KB
\end{minted}

\subsection{자동화 --- \texttt{build.bat} + \texttt{build.ps1}}

수동 5단계 절차(\texttt{Document/FinalBuild.md \S 4})를 반복하지 않도록
프로젝트 루트의 \texttt{build.bat} 한 번 더블클릭으로 \texttt{Build/}
폴더 전체가 재생산되도록 자동화했다.

\begin{minted}{text}
build.bat        (cmd 진입점, ASCII 메시지로 한국어 CP949 콘솔 호환)
  └── build.ps1  (PowerShell, 모든 빌드 로직)
        1. Build\ 폴더 초기화 + Build\vendor 생성
        2. npx --yes esbuild Web/main.js → Build/main.bundle.js
        3. Invoke-WebRequest로 Blockly 3개 다운로드
        4. dashboard.html, styles.css 복사
        5. main.html → index.html, 두 곳 패치 + Select-String 검증
\end{minted}

마지막 단계의 자체 검증은 패치 결과에 \texttt{main.bundle.js}가 포함되고
\texttt{type="module"}이 사라졌는지를 확인한다. 어느 하나라도 빠지면
\texttt{throw}로 즉시 실패하여 조용히 깨진 산출물을 만들지 않는다.

\subsection{진행 중 발견된 네 가지 함정}

오프라인 빌드는 표면적으로는 ``파일 몇 개 복사''로 보이지만 보안
컨텍스트, CORS, 인코딩, 셸 escape의 미세한 차이가 차례로 드러났다.

\begin{table}[h]
\centering
\small
\begin{tabularx}{\linewidth}{@{}lXX@{}}
\toprule
\# & 증상 & 해법 \\
\midrule
1 & \texttt{type="module"} 제거 후 일반 \texttt{<script>}가 head에서 즉시 실행되어, \texttt{elements.js}의 \texttt{document.getElementById(...)}가 DOM 준비 전 호출됨. 모든 element 참조가 \texttt{null}이라 ``아레스 탐사선 찾기'' 클릭 핸들러가 부착 못 됨. & \texttt{<script src="main.bundle.js" defer>}로 deferred 실행 보장. \\
\midrule
2 & 첫 가이드에서 \texttt{styles.css} 복사가 누락. 시각적으로 스타일이 빠진 상태로 표시. & \texttt{FinalBuild.md \S 4.4}에 명시 + \texttt{build.ps1}에 복사 단계 추가. \\
\midrule
3 & 첫 \texttt{build.bat} 시도가 PowerShell \texttt{-replace}를 cmd inline으로 호출. cmd의 \texttt{""} escape이 PowerShell에 두 글자 \texttt{""}로 전달되어 single-quote 정규식이 매치 실패. 두 번째 치환이 조용히 no-op. & 로직 전체를 별도 \texttt{build.ps1}로 분리하여 cmd-측 escape을 회피. \\
\midrule
4 & \texttt{build.ps1}의 \texttt{Get-Content -Raw}가 PowerShell 5.1 기본 인코딩(한국 Windows의 CP949)으로 \texttt{main.html}을 읽어 모든 한글이 mojibake. UTF-8로 다시 쓰지만 이미 손상된 후. & \texttt{[System.IO.File]::ReadAllText(\$path, [System.Text.Encoding]::UTF8)}로 명시적 UTF-8 읽기. PowerShell 버전·시스템 코드페이지 무관하게 동작. \\
\bottomrule
\end{tabularx}
\caption{오프라인 빌드 정착 과정에서 마주친 네 가지 함정.}
\end{table}

\begin{keypoint}
함정 1, 3, 4는 모두 ``성공적인 빌드처럼 보였지만 동작은 일부만 깨진''
패턴이다. 빌드 스크립트의 마지막에 둔 \texttt{Select-String} 자체
검증과, 시스템 코드페이지에 의존하지 않는 명시적 UTF-8 입출력이 향후
환경 변화에도 같은 함정을 다시 밟지 않도록 막아준다.
\end{keypoint}

\subsection{학습 및 운영 시사점}

\texttt{ks01.html}이 모든 코드를 인라인으로 둔 단일 파일이라 \texttt{file://}
에서 동작했던 반면, 본 빌드는 \emph{모듈 단위 코드베이스를 유지하면서도
같은 동작 표면}을 갖춘다. 학생 PC 배포 시 \texttt{Build/} 폴더만 복사하면
되고, 코드 수정 후에는 \texttt{build.bat} 한 번으로 재빌드된다.
\texttt{Document/FinalBuild.md}는 매뉴얼 절차와 자동화 양쪽을 모두
설명한다.

\section{2026-05-21 --- Web/ 12차시 × 4미션 셸 재편}

학습자 흐름을 1분기 \emph{12차시 × 4미션 = 48미션} 구조에 맞춰
\texttt{Web/} 디렉터리를 전면 재편하였다. 단일 진입(\texttt{main.html})
+ 자산 임베드를 가정하던 종전 형태에서, 랜딩 페이지(\texttt{index.html})
$\to$ 관제실(\texttt{main.html})의 2단 구성과 차시별 데이터 분리 구조로
바뀌었다. KoreaScience 데모 페이지들(\texttt{ks01.html} 외)은 본
서비스 코드와 학습 셸이 한 디렉터리에 섞이지 않도록 \texttt{KSWeb/}로
분리하였다.

\subsection{새 디렉터리 구조}

\begin{minted}{text}
Web/
├── index.html                  랜딩: "탐사선 연결" 버튼 -> main.html
├── index.css
├── main.html                   관제실: 차시/미션 네비, Blockly, dashboard iframe
├── main.js                     ES module 진입 (해시 라우터, lesson 로더 등)
├── blocklyconfig.js, commandexecutor.js, bluetooth.js,
│   constants.js, state.js, elements.js, logger.js
├── styles.css
├── overview.html               개요 탭 콘텐츠 (런타임 fetch)
├── dashboard.html              iframe 자식 (ES module 미사용)
├── Lesson01/lesson.json        1차시 4미션 데이터
├── Lesson02/lesson.json
│   ⋮
├── Lesson12/lesson.json
└── examples/
    ├── countdown.xml
    ├── greeting.xml
    └── radar_demo.xml
\end{minted}

\subsection{해시 라우터와 3-state 뷰}

\texttt{main.html}은 하나의 페이지에 세 가지 뷰(\texttt{overview} /
\texttt{lesson} / \texttt{mission})를 두고 URL hash로 전환한다.

\begin{minted}{text}
#                       -> overview (개요)
#lesson=3               -> 3차시 소개
#lesson=3&mission=2     -> 3차시 미션 2 (Blockly 워크스페이스 노출)
\end{minted}

미션 뷰 진입 시점에 한 번만 \texttt{Blockly.inject} 하고 뷰 전환은
\texttt{display}만 토글한다. ``명령 실행/저장/불러오기'' 버튼은 미션
뷰에서만 활성화된다.

\subsection{런타임 \texttt{fetch()} 의 도입 --- 새로운 외부 자원 세 종류}

이전 단일 파일 구조와 달리 컨텐츠가 외부 자원으로 분리되었기 때문에,
\texttt{main.js}는 다음 셋을 런타임에 가져온다.

\begin{table}[h]
\centering
\small
\begin{tabularx}{\linewidth}{@{}lXl@{}}
\toprule
대상 & 시점 / 용도 & 호출 위치 (main.js) \\
\midrule
\texttt{overview.html}              & 개요 뷰 최초 진입 시 1회, \texttt{dataset.loaded} 캐시 & line 249 \\
\texttt{Lesson\{NN\}/lesson.json}   & 차시/미션 진입 시 12회까지, \texttt{lessonCache} 캐시 & line 386 (\texttt{loadLesson}) \\
\texttt{examples/\{name\}.xml}      & 예제 드롭다운 선택 시 매번 & line 759 \\
\bottomrule
\end{tabularx}
\caption{재편 이후 런타임 \texttt{fetch()} 대상.}
\end{table}

\begin{keypoint}
재편 자체는 학습 흐름을 명료하게 만들었지만, 부수 효과로 ``\texttt{Build/}
한 폴더 자족 배포'' 가정이 깨졌다. 외부 자원 fetch 가 \texttt{file://}
컨텍스트에서 같은-출처 정책으로 차단되기 때문이다. 이 함정은 5월 22일
인라인 자산 shim 으로 해소한다(다음 절).
\end{keypoint}

\subsection{KSWeb/ 분리}

학습 셸과 무관한 KoreaScience 데모/소개 페이지는 \texttt{KSWeb/}로
빼냈다. \texttt{Web/}은 학습용 서비스 코드만 담는 단일 책임 디렉터리가
되었고, 빌드 자동화도 \texttt{Web/} 한 곳만 입력으로 받는다.

\section{2026-05-22 --- 오프라인 빌드의 신구조 대응: 인라인 자산 shim}

5월 20일 정착한 \texttt{Build/} 자족 배포본은 5월 21일 재편 이후 더
이상 실제 서비스를 반영하지 못했다. \texttt{Web/main.js}가 추가로
\texttt{fetch()} 하게 된 세 자원이 \texttt{file://}에서 차단되어, 빌드는
``성공''하지만 학생 PC에서 차시를 한 번도 펼칠 수 없는 상태가 된다. 본
절은 \texttt{Web/} 소스를 일절 수정하지 않은 채 \texttt{file://}에서도
모든 \texttt{fetch()} 가 성공하도록 만든 \emph{인라인 자산 shim} 빌드
전략을 정리한다.

\subsection{왜 단순 esbuild 재번들로 끝나지 않는가}

이전 빌드는 \texttt{main.html} 단일 진입 + 자산 임베드를 가정했다.
재편된 \texttt{main.js}가 fetch 하는 세 종류는 \texttt{import} 그래프에
들어가 있지 않기 때문에 \texttt{esbuild --bundle}로 잡히지 않는다.
\texttt{file://}에서 \texttt{fetch()}는 \texttt{origin: null}의 같은-출처
정책으로 무조건 차단된다.

\begin{table}[h]
\centering
\small
\begin{tabularx}{\linewidth}{@{}lXX@{}}
\toprule
의존성 & \texttt{file://}에서의 문제 & 본 빌드의 해법 \\
\midrule
\texttt{Web/} 8개 ES module               & \texttt{<script type="module">} 같은-출처 차단 & esbuild IIFE 번들 (\texttt{main.bundle.js}) \emph{(2026-05-20)} \\
Blockly CDN 3개 스크립트                   & 인터넷 필요 & \texttt{Build/vendor/} 사본, 상대 경로 참조 \emph{(2026-05-20)} \\
\texttt{overview.html} / \texttt{Lesson\{NN\}/lesson.json} / \texttt{examples/*.xml}
  & 런타임 \texttt{fetch()} 가 같은-출처로 차단 (\texttt{Web/} 8 module 차단과 같은 원인, 다른 형태)
  & 빌드 시점 인라인 + \texttt{window.fetch} 가로채기 shim \emph{(본 절)} \\
\bottomrule
\end{tabularx}
\caption{\texttt{file://} 동작을 막는 의존성 세 묶음과 시기별 해법.}
\end{table}

\subsection{해법 --- \texttt{vendor/inline\_assets.js} shim}

빌드 시점에 모든 \texttt{fetch()} 대상을 한 JS 파일에 문자열로 인라인하고,
런타임에 \texttt{window.fetch} 를 가로채는 얇은 shim 을 끼워 넣는다.
번들 코드는 일절 수정하지 않으며, \texttt{Web/} 소스 트리도 그대로다.

\begin{minted}{javascript}
// Build/vendor/inline_assets.js (build.ps1이 자동 생성)
(function () {
  if (window.__ARES_INLINE_INSTALLED__) return;
  window.__ARES_INLINE_INSTALLED__ = true;
  var DATA = {};
  DATA['overview.html'] = "...";
  DATA['Lesson01/lesson.json'] = "{...}";
  /*  ... 12 lessons + 3 examples ... */

  var origFetch = window.fetch ? window.fetch.bind(window) : null;
  function lookup(input) {
    var url = (typeof input === 'string') ? input : (input && input.url) || '';
    url = String(url).split('?')[0].split('#')[0];
    for (var key in DATA) {
      if (!Object.prototype.hasOwnProperty.call(DATA, key)) continue;
      if (url === key) return key;
      if (url.endsWith('/' + key)) return key;
      if (url.endsWith(key)) return key;
    }
    return null;
  }
  function mimeFor(key) {
    if (key.endsWith('.json')) return 'application/json';
    if (key.endsWith('.xml'))  return 'application/xml';
    return 'text/html; charset=utf-8';
  }
  window.fetch = function (input, init) {
    var key = lookup(input);
    if (key !== null) {
      return Promise.resolve(new Response(DATA[key], {
        status: 200,
        headers: { 'Content-Type': mimeFor(key) }
      }));
    }
    if (origFetch) return origFetch(input, init);
    return Promise.reject(new Error('fetch unsupported (no inline match)'));
  };
})();
\end{minted}

\texttt{main.html}의 head에 다음 두 줄을 \texttt{defer} 순서로 끼워
\texttt{main.bundle.js} \emph{이전}에 shim 이 설치되도록 보장한다.

\begin{minted}{html}
<script src="vendor/inline_assets.js" defer></script>
<script src="main.bundle.js" defer></script>
\end{minted}

\begin{keypoint}
\texttt{defer} 속성을 갖는 외부 스크립트는 \emph{문서 내 등장 순서대로}
실행되므로 두 줄의 순서가 곧 실행 순서다. shim이 먼저 실행되어
\texttt{window.fetch}를 교체한 뒤 번들이 진입하므로, 번들 안에서 일어나는
모든 \texttt{fetch()} 호출이 인라인 데이터를 본다.
\end{keypoint}

\subsection{경로 매칭 규칙 --- \texttt{endsWith} 세 단계}

\texttt{main.js}의 fetch 호출은 세 가지 모양으로 일어난다.

\begin{itemize}
  \item \texttt{fetch('overview.html', ...)} --- 순수 상대 경로
  \item \texttt{fetch(`Lesson\$\{pad\}/lesson.json`, ...)} --- 하위 폴더 포함 상대 경로
  \item \texttt{fetch(new URL(`examples/\$\{name\}.xml`, window.location.href).href)} --- 절대 URL
\end{itemize}

\texttt{file://} 컨텍스트에서는 모두 \texttt{file:///.../Build/<key>} 형태의
절대 URL로 정규화된다. shim의 \texttt{lookup()}은 우선 query/hash 를
잘라낸 뒤 \texttt{url === key}, \texttt{url.endsWith('/' + key)},
\texttt{url.endsWith(key)} 순으로 시도한다. 첫 번째는 키와 동일한 짧은
입력 보호용, 두 번째는 디렉터리 경계 보장(\texttt{lesson.json} 키가
\texttt{xlesson.json} URL에 매치되지 않도록), 세 번째는 키가 절대 URL의
끝부분과 정확히 같을 때를 위한 안전망이다.

\subsection{문자열 인라인 시의 두 가지 작은 함정}

\texttt{PowerShell}에서 한국어 텍스트를 JS 문자열 리터럴로 안전하게 옮기는
과정에서 다음 두 함정이 있었다.

\begin{table}[h]
\centering
\small
\begin{tabularx}{\linewidth}{@{}lXX@{}}
\toprule
\# & 증상 & 해법 \\
\midrule
1 & 자료(예: \texttt{overview.html}) 안의 \texttt{<\/...}{}>가 \texttt{<script>} 안에서 그대로 출력되면 HTML 파서가 inline\_assets.js 스크립트를 \emph{거기서 종료}시켜 그 뒤가 평문으로 새어 나옴. & \texttt{ConvertTo-Json} 출력에 \texttt{Replace('<\/', '<\textbackslash{}/') } 후처리. JS 리터럴 안의 \texttt{"<\textbackslash{}/script>"}는 동일한 문자열이지만 HTML 파서에 끝-태그로 보이지 않는다. \\
\midrule
2 & PowerShell의 \texttt{Set-Location \$PSScriptRoot}이 PSDrive cwd만 바꾸고 .NET 의 \texttt{Environment.CurrentDirectory}는 그대로 둠. \texttt{[System.IO.File]::WriteAllText('Build\textbackslash{}index.html', ...)}처럼 상대 경로를 쓰면 호출자의 셸 cwd(예: \texttt{Missions\textbackslash{}})에서 풀려 \texttt{DirectoryNotFoundException}. & \texttt{Set-Location} 직후 \texttt{[System.Environment]::CurrentDirectory = \$PSScriptRoot} 동기화. \\
\bottomrule
\end{tabularx}
\caption{인라인 자산 생성 단계의 두 함정.}
\end{table}

\subsection{새 \texttt{Build/} 구조}

빌드 결과는 \emph{서비스의 2단 진입을 그대로 보존}한다 ---
\texttt{Build/index.html}(랜딩) $\to$ \texttt{Build/main.html}(관제실).
이전 빌드가 \texttt{main.html}을 \texttt{Build/index.html}로 \emph{대체}했던
구조와 다르다. 랜딩의 \texttt{<a href="../index.html">프로젝트 소개로
돌아가기</a>} 백링크는 \texttt{Build/} 단독 배포에서 갈 곳이 없으므로
빌드 패치에서 제거한다.

\begin{minted}{text}
Build/                          ~1.2 MB (self-contained)
├── index.html                  랜딩 (Web/index.html 패치)
├── main.html                   관제실 (Web/main.html 패치:
│                                       Blockly vendor/ 경로 +
│                                       inline_assets.js + main.bundle.js, 둘 다 defer)
├── main.bundle.js              92 KB  (8개 ES module IIFE 번들)
├── dashboard.html              (수정 없음, iframe 자식)
├── styles.css, index.css
└── vendor/
    ├── inline_assets.js        100 KB (overview + 12 lesson.json + 3 examples + fetch shim)
    ├── blockly_compressed.js
    ├── blocks_compressed.js
    └── python_compressed.js
\end{minted}

\subsection{\texttt{build.ps1} 의 7단계}

\begin{minted}{text}
build.bat              (cmd 진입점, 메시지 갱신)
  └── build.ps1
        1. Build\ 폴더 초기화 + Build\vendor 생성
        2. npx --yes esbuild Web/main.js -> Build/main.bundle.js  (IIFE, es2018)
        3. Invoke-WebRequest 로 Blockly 3개 다운로드 -> Build/vendor/
        4. dashboard.html, styles.css, index.css 복사
        5. Web/index.html -> Build/index.html (랜딩, "../index.html" 백링크 제거)
        6. Web/main.html  -> Build/main.html
           - Blockly CDN URL -> vendor/ 상대 경로
           - <script type="module" src="main.js"> ->
             <script src="vendor/inline_assets.js" defer></script>
             <script src="main.bundle.js" defer></script>
           - Select-String 자체 검증 3종 (main.bundle.js 포함, type="module" 부재, inline_assets.js 포함)
        7. vendor/inline_assets.js 생성
           - overview.html 1개 + Lesson{01..12}/lesson.json 12개 + examples/*.xml N개
           - ConvertTo-Json 으로 JS 리터럴 escape + "</" -> "<\/" 후처리
           - 항목 수 sanity check (>= 13)
\end{minted}

\subsection{검증}

\begin{enumerate}
  \item \texttt{Build/} 폴더를 USB 등으로 인터넷 없는 PC에 복사.
  \item \texttt{Build/index.html}을 Chrome 또는 Edge에서 더블클릭.
        URL 표시줄이 \texttt{file:///.../Build/index.html}로 시작하는지 확인.
  \item ``탐사선 연결'' 버튼 클릭 $\to$ \texttt{Build/main.html} 진입.
  \item 콘솔(F12)에 CORS / fetch 실패가 없는지 확인. ``개요'' 화면이
        정상 표시되어야 한다(\texttt{overview.html} 인라인 동작 확인).
  \item 차시 드롭다운에서 1차시 선택 $\to$ 4개 미션 목록 표시.
        (\texttt{Lesson01/lesson.json} 인라인 동작 확인).
  \item 미션 1 진입 $\to$ Blockly 워크스페이스 노출. 예제 드롭다운에서
        ``카운트다운'' 선택 $\to$ 블록 로드. (\texttt{examples/countdown.xml}
        인라인 동작 확인).
  \item ``신호 연결'' 클릭 $\to$ BT05/HM-10 페어링 $\to$ ``명령 전송''으로
        LED 블록 1개 실행.
  \item ``한꺼번에 실행'' 블록으로 카운트다운 시퀀스 실행 $\to$ BATCH
        경로 정상 확인.
\end{enumerate}

\subsection{개발자 운영 안내}

\begin{itemize}
  \item \textbf{개발 시}는 \texttt{Web/}을 \texttt{python -m http.server 8000}
        등으로 직접 서빙. \texttt{Build/}는 \emph{빌드 산출물}이므로
        편집 대상이 아니다.
  \item \texttt{Web/Lesson\{NN\}/lesson.json}이나 \texttt{Web/examples/*.xml}
        을 수정하면 반드시 \texttt{build.bat} 재실행. \texttt{Build/}
        쪽의 인라인 데이터는 빌드 시점 snapshot 이기 때문이다.
  \item \texttt{Web/} 트리에 새 차시 데이터(\texttt{LessonNN/lesson.json})
        를 더해도 \texttt{build.ps1}은 1\,$\sim$\,12까지 자동 순회하므로
        파일만 추가하면 된다. \texttt{LESSON\_CATALOG}(main.js)와 일관
        유지만 챙기면 됨.
\end{itemize}

\section{변경 파일 요약}

\begin{table}[h]
\centering
\small
\begin{tabularx}{\linewidth}{@{}llX@{}}
\toprule
일자 & 파일 & 변경 요지 \\
\midrule
05-19 & \texttt{Web/commandexecutor.js} & \texttt{FIRE\_AND\_FORGET\_HEADS} 집합 도입, \texttt{sendCommand()}에서 분기 처리, cooldown 분리(40 / 20\,ms). \\
05-19 & \texttt{Pico/main.py} & \texttt{NO\_RESPONSE\_CMDS} 클래스 변수, \texttt{\_read\_uart\_line()} 폴링 방식으로 교체, \texttt{\_needs\_response()} 도입, \texttt{run()}에서 조건부 \texttt{\_send\_response()}. \\
\midrule
05-20 & \texttt{Pico/main.py} & \texttt{boot()}의 OLED 호출 None 가드 추가. 05-19 BLE 최적화 재적용. \\
05-20 & \texttt{Pico/process\_data.py} & 누락되어 있던 \texttt{\_handle\_sys\_set()} 메서드 헤더 복구. \\
05-20 & \texttt{Pico/dcmotor.py} & \texttt{DIR+PWM} 단일 변조에서 양 핀 모두 PWM으로 전환. \texttt{dc\_forward}/\texttt{dc\_backward}가 한쪽 핀을 \texttt{PWM\_MAX} 고정, 다른 쪽을 \texttt{PWM\_MAX-speed}로 변조하여 양방향 brake-modulated 대칭 출력. \\
\midrule
05-20 & \texttt{Web/blocklyconfig.js} & \texttt{batch\_block} JSON 정의 + \texttt{BATCH\_FORBIDDEN\_TYPES} + \texttt{attachBatchBlockValidator()}. \\
05-20 & \texttt{Web/main.js} & \texttt{attachBatchBlockValidator(Blockly)} 호출 등록. \\
05-20 & \texttt{Web/commandexecutor.js} & \texttt{processBlock}에서 \texttt{batch\_block} 분기, \texttt{\_processBatch} 헬퍼 (자식 평탄화 + 금지 블록 차단 + 단일 \texttt{BATCH;...} 송신). \\
05-20 & \texttt{Web/main.html} & toolbox에 ``🚀 한꺼번에'' 카테고리 추가. \\
05-20 & \texttt{Web/constants.js} & \texttt{CHUNK\_DELAY} 50 $\rightarrow$ 100\,ms. \\
05-20 & \texttt{Web/bluetooth.js} & \texttt{writeValueWithoutResponse} 우선 사용(fallback 포함), \texttt{pendingResolve}를 송신 \emph{전} 설치(race 해소), 마지막 청크 후 \texttt{CHUNK\_DELAY} guard. \\
05-20 & \texttt{Pico/process\_data.py} & \texttt{BATCH;} 디스패처 + \texttt{\_handle\_batch()} (파이프 분리 후 기존 \texttt{process()} 재진입). \\
05-20 & \texttt{Pico/main.py} & \texttt{RECEIVE\_TIMEOUT\_MS} 50 $\rightarrow$ 500\,ms, 타임아웃 시 부분 데이터 반환 제거(buffer 유지). \\
\midrule
05-20 & \texttt{build.bat} (root) & cmd 진입점. \texttt{build.ps1} 호출 후 결과 \texttt{pause}. ASCII 메시지로 CP949 콘솔 호환. \\
05-20 & \texttt{build.ps1} (root) & 빌드 5단계 + 자체 검증(\texttt{Select-String}). UTF-8 명시 입출력. \\
05-20 & \texttt{Document/FinalBuild.md} & 오프라인 (\texttt{file://}) 배포 빌드 가이드. 매뉴얼 절차와 자동화 안내. \\
05-20 & \texttt{Build/} & 자족 배포 산출물(\texttt{index.html}, \texttt{main.bundle.js}, \texttt{styles.css}, \texttt{dashboard.html}, \texttt{vendor/*.js}). \\
\midrule
05-21 & \texttt{Web/index.html} (신규) & 랜딩 페이지. ``탐사선 연결'' 버튼으로 \texttt{main.html} 진입. \\
05-21 & \texttt{Web/index.css} (신규) & 랜딩 전용 스타일. \\
05-21 & \texttt{Web/main.html} & 차시/미션 네비, 3-state 뷰(\texttt{overview}/\texttt{lesson}/\texttt{mission}), iframe \texttt{dashboard.html} 임베드. \\
05-21 & \texttt{Web/main.js} & 해시 라우터(\texttt{\#lesson=N\&mission=M}), \texttt{LESSON\_CATALOG}, \texttt{loadLesson()} 캐시, \texttt{enterOverview/Lesson/Mission}. 신규 \texttt{fetch()} 3종 도입. \\
05-21 & \texttt{Web/Lesson\{01..12\}/lesson.json} (신규) & 차시별 4미션 데이터(스토리, 학습 목표, 샘플 코드, 요약). \\
05-21 & \texttt{Web/overview.html} (신규) & 개요 뷰 콘텐츠. \\
05-21 & \texttt{Web/examples/*.xml} & \texttt{greeting}, \texttt{countdown}, \texttt{radar\_demo} 등 예제 블록 XML. \\
05-21 & \texttt{KSWeb/} (신규) & KoreaScience 데모/소개 페이지 분리 보관. \\
05-21 & \texttt{MISSIONS.md} (신규) & 1분기 12차시 × 4미션 = 48미션 커리큘럼 요약. \\
\midrule
05-22 & \texttt{build.ps1} (root) & 5단계 $\to$ 7단계 확장. \texttt{Web/index.html} $\to$ \texttt{Build/index.html}(랜딩, 백링크 제거), \texttt{Web/main.html} $\to$ \texttt{Build/main.html} 패치(Blockly vendor + \texttt{inline\_assets.js} + \texttt{main.bundle.js}, 둘 다 defer), 7단계로 \texttt{vendor/inline\_assets.js} 생성(overview + 12 lesson + 예제 + fetch shim). \texttt{[Environment]::CurrentDirectory} 동기화로 셸 cwd 비의존. \texttt{ConvertTo-Json} 결과의 \texttt{<\/} $\to$ \texttt{<\textbackslash{}/} 후처리. \\
05-22 & \texttt{build.bat} (root) & 헤더 주석을 새 산출물 구조(랜딩 \texttt{index.html} + 관제실 \texttt{main.html} + \texttt{vendor/inline\_assets.js})에 맞춰 갱신. 빌드 성공 메시지 변경. \\
05-22 & \texttt{Build/index.html}, \texttt{Build/main.html}, \texttt{Build/vendor/inline\_assets.js} & 새 자족 배포본의 진입/관제실/인라인 shim 산출물. \\
\midrule
05-28 & \texttt{Web/simulation.js} (신규) & \texttt{Web/main.js} 의 3D 시뮬레이션 파트(약 520줄: \texttt{TOPICS}, \texttt{buildSim}, \texttt{setupSimulation})를 ES 모듈로 분리. \texttt{CommandExecutor} 만 직접 import 하고 \texttt{workspace} 는 인자로 주입. \\
05-28 & \texttt{Web/main.js} & 시뮬 코드 제거 후 \texttt{setupSimulation(\{ workspace \})} 호출과 \texttt{simController} 보관만 유지(1455 $\to$ 937 라인, $-$518). \\
05-28 & \texttt{Web/main.html} & 시뮬 카드 헤더에 \texttt{.sim-launch-buttons} 그룹 추가 — \texttt{\#simRadar}(레이더 토글), \texttt{\#simRocket}(로켓 발사 토글). \\
05-28 & \texttt{Web/styles.css} & \texttt{.sim-launch-btn} 외관 등록(우주 신호등 버튼과 동일 형태, 활성색만 청록 \texttt{\#3cbfbf}). \\
05-28 & \texttt{Web/simulation.js} (LaunchStation 후처리) & 단일 mesh \texttt{LaunchStation.glb} 에서 bbox 휴리스틱으로 안테나(fx 0.78\textendash 0.92, fy>0.70)와 로켓(fx 0.28\textendash 0.46, fy>0.68)을 정점 단위로 분리. 안테나: 회색 + y축 \texttt{rotation} pivot(centroid - 0.01 x 보정). 로켓: 노란색 + 하단 화염 sprite(additive) + \texttt{PointLight}. \\
05-28 & \texttt{Web/simulation.js} (애니메이션) & 레이더 가동 시 \texttt{rotation.y += 0.15}(${\approx}500^{\circ}$/s). 로켓 발사 시 \texttt{rocketGroup.position.y} 를 \texttt{ROCKET\_RISE=10} 까지 비대칭 ease 로 이동 — 발사 ease-out, 중지 ease-in(t$^2$). 카메라 \texttt{controls.target} 을 로켓 world centroid 로 직접 락온하고 \texttt{camera.position.y} 도 동일 상승. 중지 후 saved 상태로 스냅 복귀. 속도 \texttt{ROCKET\_SPEED=0.00267}(${\approx}6$초 완주). \\
05-28 & \texttt{build.ps1} & 시뮬 GLB 7종(\texttt{AlbiStaticLow}, \texttt{LampBox}, \texttt{LampGeneral}, \texttt{LampHand1\textendash 3}, \texttt{LaunchStation}) base64 BIN 일괄 인라인. \texttt{Build/Mesh/} 폴더 폐기, \texttt{ares\_robot.embed.js} 를 \texttt{Build/} 루트로 복사하고 \texttt{ares\_robot.glb} fallback 제거. 랜딩 \texttt{index.html} 의 \texttt{Mesh/ares\_robot.*} 경로를 루트로 치환. \\
05-28 & \texttt{Document/FinalBuild.md} & 폴더 트리에서 \texttt{Mesh/} 항목 제거, \texttt{ares\_robot.embed.js} 를 루트로 옮기고 시뮬 GLB 가 fetch shim 으로 처리됨을 명시. \\
05-28 & \texttt{Document/CodeChangeHistory.tex} & 05-28 커밋 해시(\texttt{110acf1}, \texttt{9e04b2b})를 반영해 (pending) 행 치환. ``우선 할 일 (Backlog)'' 섹션 신설 — \texttt{Build/vendor/inline\_assets.js} 50.54\,MB(GitHub 권장 한도 50\,MB 초과)에 대한 GLB 최적화 항목 등록(Meshopt/Draco 압축, 텍스처 축소, \texttt{Build/} \texttt{.gitignore} 처리 옵션 등). \\
\midrule
06-02 & \texttt{Web/simulation.js} & 알비 가슴 LED(붉은 \texttt{CHEST\_PALETTE} 구) 1개 추가. \texttt{applyTopicEffect} 에서 시뮬 명령 $\to$ 시각 반응 매핑: \texttt{LED\_ON}/\texttt{LED\_OFF}/LED 패턴 $\to$ 눈·슬롯·발사대 LED, \texttt{BUZZER\_ON} $\to$ 가슴 LED 점등 + 부저 비프(\texttt{playBeep}), \texttt{SLEEP} $\to$ hold 대기. \\
\midrule
06-04 & \texttt{Web/simulation.js} & 발사대 LED 6개 동적 생성 --- LED1\textendash 5(전면 세로 줄, 초록 \texttt{LAUNCH\_STRIP\_PALETTE}) + LED0(로켓 바닥 도넛, 빨강 \texttt{LAUNCH\_TORUS\_PALETTE}). \texttt{applyLed} 채도 튜닝 확장: \texttt{intensityScale}(발광량 $\downarrow$)·\texttt{opacityOn}(구체 불투명도 $\uparrow$)·\texttt{glowScale}(가산 글로우 $\downarrow$)로 ACES 흰빛 날림 억제. LED0 을 가슴 팔레트에서 전용 팔레트로 분리. \\
06-04 & \texttt{Web/main.html}, \texttt{Web/styles.css} & 발사대 LED 전체 토글 버튼(\texttt{\#simLaunchLeds}, \texttt{.sim-launch-led-btn}) --- 한 번에 LED0\textendash 5 점등/소등, 활성색 발광 빨강. \\
06-04 & \texttt{Web/simulation.js} & 로켓 발사 시 회백색 연기/구름 분출 --- Sprite 풀 80개 + 뭉게구름 텍스처(다중 radial blob). \texttt{rocketMeshRef}(mesh-local)에 부착해 로켓과 함께 솟지 않고 발사대 바닥 구름 + 상승 trail 형성. 추진 중 분출 후 팽창·페이드·소멸. \\
06-04 & \texttt{Web/simulation.js} & 시뮬 실행 종료(\texttt{finally})·시뮬 카드 닫기(\texttt{close}) 시 떠 있는 로켓을 ``발사 중지''와 동일하게 원위치 복귀시키는 과정 재생. \texttt{rocketAtRest} getter + \texttt{finalizeClose}(복귀 완료 후 카드 숨김). \\
06-04 & \texttt{Web/simulation.js} & 로켓 발사음 \texttt{playRocketLaunch()} --- Web Audio 화이트노이즈를 lowpass(럼블)·bandpass(분사 로어)로 필터링한 ${\approx}3.6$초 절차적 굉음. UI 로켓 버튼·\texttt{GUN\_FIRE} 명령에서 발사 시작 시 재생. \\
06-04 & \texttt{Web/examples/albi\_song.xml} (신규) & 예제 ``알비의 노래'' --- Real Madrid 응원가 \textit{Hala Madrid y nada más} 1절 주선율을 \texttt{buzzer\_on} 40음으로 전사(반음은 정확 Hz, 예: C\#5=554). \\
06-04 & \texttt{Web/main.html} & 예제 드롭다운(\texttt{\#exampleSelect})에 ``🎵 알비의 노래'' 옵션 추가. \\
\bottomrule
\end{tabularx}
\caption{2026-05-19 \textendash{} 06-04 변경 범위 (롤링).}
\end{table}

\begin{table}[h]
\centering
\small
\begin{tabularx}{\linewidth}{@{}llX@{}}
\toprule
일자 & 커밋 & 메시지 \\
\midrule
05-19 & \texttt{7e5b93f} & Optimize BLE pipeline: fire-and-forget for output commands \\
05-19 & \texttt{110bb27} & Revert ``Optimize BLE pipeline: fire-and-forget for output commands'' \\
05-20 & \texttt{600243a} & Re-apply BLE pipeline optimization with OLED boot crash fix \\
05-20 & \texttt{de9cffc} & Fix DC motor: forward PWM asymmetry and missing SYS\_SET handler \\
05-20 & \texttt{5780485} & Drive H-bridge with dual PWM for symmetric DC motor output \\
05-20 & \texttt{4fccfe7} & Add ``한꺼번에 실행'' (BATCH) block to bundle commands into one BLE write \\
05-20 & \texttt{c6d57ae} & Hold UART buffer across timeouts so multi-chunk BATCH stays intact \\
05-20 & \texttt{8d341af} & Settle BATCH path: drop debug prints, switch BLE writes to without-response \\
05-20 & \texttt{21c4e26} & Fix BLE response race: install pendingResolve before sending \\
05-20 & \texttt{233f185} & Skip CHUNK\_DELAY after the final chunk so single-chunk commands stay fast \\
05-20 & \texttt{328c81a} & Add FinalBuild.md --- offline (file://) packaging guide \\
05-20 & \texttt{eea8b54} & Produce Build/ offline artifact and patch FinalBuild.md for styles.css \\
05-20 & \texttt{e6a35ce} & Add defer to Build/index.html bundle <script> so DOM is ready first \\
05-20 & \texttt{0af8ac2} & Add build.bat / build.ps1 to automate the FinalBuild procedure \\
05-20 & \texttt{316aa2b} & Read main.html as UTF-8 in build.ps1 so Korean text survives the patch \\
05-20 & \texttt{ff46f92} & Add FinalBuild section to CodeChange20260520 as the last work chapter \\
\midrule
05-21 & \texttt{0535c89} & Move KoreaScience-specific pages out of Web/ into new KSWeb/ \\
05-21 & \texttt{c1b22df} & Add MISSIONS.md --- 1분기 12차시 × 4미션 커리큘럼 요약 \\
05-21 & \texttt{d3d6785} & Restructure Web/ as 12-lesson mission-based service \\
\midrule
05-22 & \textit{(pending)} & Rebuild offline package for restructured Web/: inline\_assets.js fetch shim, two-stage entry (index.html landing $\to$ main.html), 7-step build.ps1, build.bat header refresh \\
05-22 & \textit{(pending)} & Rename CodeChange20260520 $\to$ CodeChangeHistory and append 05-21 / 05-22 chapters \\
\midrule
05-28 & \textit{(pending)} & Split simulation out of Web/main.js into Web/simulation.js and bundle launchpad scene (antenna recolor + radar spin, rocket recolor + launch animation with camera tracking and flame) \\
05-28 & \texttt{110acf1} & Sim: simulation.js 분리 + 발사대 안테나/로켓 액션 추가 \\
05-28 & \texttt{9e04b2b} & Build: 시뮬 GLB 7종 inline\_assets 통합, Build/Mesh 폴더 폐기 \\
\midrule
06-02 & \texttt{4670c8a} & Sim: 알비 가슴 LED + LED\_ON/OFF/BUZZER\_ON/SLEEP 시뮬 명령 반응 \\
\midrule
06-04 & \texttt{d26629f} & Sim: 발사대 LED 전체 토글 버튼 추가 (LED0 도넛 + LED1\textendash 5 세로 줄) \\
06-04 & \texttt{72bf32b} & Sim: 로켓 발사 연기/구름 + 종료 시 로켓 자동 복귀 + 발사대 LED 채도 강화 \\
06-04 & \texttt{e666ea9} & Sim: 로켓 발사 시 절차적 발사음 추가 (럼블+분사 굉음) \\
06-04 & \texttt{db9985b} & Example: ``알비의 노래'' 추가 --- 부저로 Hala Madrid(알라 마드리드) 응원가 후렴 연주 \\
06-04 & \texttt{0f4fe08} & Example: ``알비의 노래''를 악보 기반 1절 주선율로 교체 (Hala Madrid y nada más) \\
\bottomrule
\end{tabularx}
\caption{관련 커밋 (\texttt{main} 브랜치). 05-22 행만 본 문서 갱신 시점에 아직 커밋되지 않은 작업.}
\end{table}

\section*{우선 할 일 (Backlog)}

다음 작업은 아직 착수 전이며, 다음 세션의 우선 처리 대상이다.

\begin{table}[h]
\centering
\small
\begin{tabularx}{\linewidth}{@{}llX@{}}
\toprule
등록일 & 항목 & 내용 \\
\midrule
05-28 & GLB 최적화 & \texttt{Build/vendor/inline\_assets.js} 가 50.54\,MB 로 GitHub 권장 한도 50\,MB 초과(푸시는 성공, 절대 한도 100\,MB 미만). 원인: 시뮬 GLB 7종 base64 인라인. 주범 \texttt{AlbiStaticLow.glb}(14.6\,MB), \texttt{LaunchStation.glb}(14.8\,MB), \texttt{LampBox.glb}(9.6\,MB). 해법 후보: (1) \texttt{gltfpack -cc -tc} (Meshopt) 또는 Draco 압축으로 GLB 자체 경량화 — 목표 합산 $\le$10\,MB, (2) 각 모델에서 카메라가 보지 않는 디테일/큰 텍스처 정리(예: 발사대 박스 내부), (3) 정책 변경으로 \texttt{Build/} 를 \texttt{.gitignore} 처리하고 GitHub Releases 로 배포(현재 추적 중인 운영을 같이 바꿔야 함). 경량화 후 \texttt{build.bat} 재실행 $\to$ \texttt{inline\_assets.js} 크기 재확인. \\
\midrule
06-04 & \texttt{Build/} 재빌드 & 06-04 \texttt{Web/} 변경(\texttt{simulation.js} 로켓 연기/발사음/복귀·발사대 LED, \texttt{examples/albi\_song.xml} 신규, \texttt{main.html} 예제 옵션)은 아직 \texttt{Build/} 오프라인 산출물에 미반영. 배포본 갱신 시 \texttt{build.bat} 재실행 필요(인라인 데이터는 빌드 시점 snapshot). \\
\bottomrule
\end{tabularx}
\caption{우선 할 일 목록 — 완료 시 위쪽 변경 행으로 이동.}
\end{table}

\end{document}
